\subsection{The remaining $81$-torsion comparison has only two inertia pairs}

\begin{lemma}[Norm and prime-to-$3$ characters in the quaternion algebra]
\label{lem:B-RD32-norm}
For the finite subgroups $H$ in Lemma~\ref{lem:B-RD31-finite},
the reduced norm is trivial whenever $3\mid |H|$.
If it is nontrivial, $H$ is cyclic of order $8$.
If it is trivial, the possible orders are
\[
                           1,2,3,4,6,12.
                                                        \tag{B-RD32.1}
\]
In the norm-one case the prime-to-$3$ character depends only on
the tame quotient order $t\in\{1,2,4\}$, and on a prime-to-$3$
element whose image in that quotient has order $r$ it is
\[
       r=1:\ 2,\qquad r=2:\ -2,\qquad r=4:\ 0.
                                                        \tag{B-RD32.2}
\]
Thus two norm-one special inertia images with the same tame
quotient order have identical prime-to-$3$ characters as
functions on inertia.
\end{lemma}
\begin{proof}
Retain the explicit quaternion algebra (B-RD31.1). A noncentral
element has characteristic polynomial of degree two, whose
constant term is its reduced norm. There is only one element
of order two, namely $-1$, because the algebra is a division
algebra and $(x-1)(x+1)=0$ forces $x=1$ or $-1$.

If the wild subgroup is nontrivial, it is $P=C_3$. Its
centralizer in $H$ is contained in
$\mu_6\subset\mathbf Q_3(\zeta_3)$, on which the reduced norm
is $1$. If $H$ is not equal to this centralizer, its quotient
by the centralizer has order two. The tame quotient of order
two cannot give a nonabelian group of order six, since its
involutions would be noncentral. Therefore the nonabelian
case has order twelve. Its Sylow $2$-subgroup injects under
reduction into $\mathbf F_9^\times$, so it is cyclic of order
four. A generator $j$ satisfies $j^2=-1$ and has reduced
characteristic polynomial $X^2+1$, hence reduced norm $1$.
The centralizer and $j$ generate $H$. This proves norm one
for every wild $H$.

If $H$ is tame, reduction embeds it in $\mathbf F_9^\times$,
so it is cyclic of order dividing eight. A generator lies in
a quadratic commutative subfield. For an element of order
eight this is the unramified quadratic field, and its reduced
norm is $\zeta_8^{1+3}=\zeta_8^4=-1$. For elements of order
dividing four the norm is $1$. This proves (B-RD32.1) and
the assertion about nontrivial norm.

In a norm-one subgroup, a prime-to-$3$ element has order
$1,2$, or $4$. Its trace is respectively $2,-2$, or $0$,
the last two facts following from $-1$ and $X^2+1$.
Its cyclic subgroup maps injectively to the tame quotient,
so its order equals the order of its image there.
Over the completed maximal unramified base, the absolute
tame inertia is procyclic and has a unique quotient of each
of the orders $1,2,4$. The order of the image of an element
in that quotient is independent of the choice of a generator.
Formula (B-RD32.2) therefore proves equality as functions on
inertia, even when the two wild good-reduction fields differ.
\end{proof}

\begin{proposition}[B-RD32: exact restriction of the two undecided tests]
\label{prop:B-RD32}
\status{Proved; parent and independent second proof audits passed.}
Keep $p=3$, $e_K=2$, $\mu_3\subset K$, and height two.
For a pair with isomorphic $27$-torsion and different
prime-to-$3$ characters, both groups must be potentially
supersingular, both finite special inertia images must have
trivial reduced norm, their tame quotient orders must differ,
and their joint faithful good-reduction image must have order
strictly greater than $9$.

For a pair with isomorphic $81$-torsion and different
prime-to-$3$ characters, after interchanging the pair their
individual faithful special inertia orders must be
\[
                             (12,3)\quad\hbox{or}\quad(12,6),
                                                        \tag{B-RD32.3}
\]
and their joint faithful special inertia image must have
order exactly $36$. In particular its wild subgroup has
order $9$ and the two wild order-three quotients are
independent.

Consequently $\nu_{\min}(3,2,2)\leq3$ if and only if
no $81$-torsion comparison of the two specified kinds has
different prime-to-$3$ characters. This is a restricted
remaining test, not a claim that either kind exists.
\end{proposition}
\begin{proof}
Theorem~\ref{thm:B-RD31} reduces either test to two whole
potentially supersingular groups. Its determinant argument
shows that their finite determinant characters agree.
Under crystalline comparison this finite determinant is
the reduced norm of the special quaternionic inertia
action: the comparison duality or period Tate twist does
not change finite descent inertia.

If this determinant is nontrivial, Lemma~\ref{lem:B-RD32-norm}
makes both special images cyclic of order eight and tame.
Their common good extension then has ramification at most
eight, hence absolute ramification at most sixteen.
Bondarko's Theorem~E(1), with $u=2$ and $\ell=1$, applies
since $16\leq2\cdot3^2$. Already an isomorphism on
$27$-torsion makes their special $3$-torsion equivariantly
isomorphic and their prime-to-$3$ characters equal.
Thus any counterexample has norm one on both sides.

Their tame quotient orders must differ by
Lemma~\ref{lem:B-RD32-norm}. Write $m$ for the order of
the joint faithful image. Its wild subgroup has order
$1,3$, or $9$, and its tame quotient has order $1,2$, or
$4$. In particular
\[
                m\in\{1,2,3,4,6,9,12,18,36\}.
                                                        \tag{B-RD32.4}
\]
For $m\leq9$, the common absolute ramification is at most
$18=2\cdot3^2$, so the same $u=2$ argument settles the
$27$-torsion comparison. This proves the first assertion.

At level $81$, take $u=3$ and $\ell=1$. The comparison is
settled whenever $2m\leq2\cdot3^3$, or $m\leq27$.
The only larger value in (B-RD32.4) is $36$. This forces
both wild images to have order three and the joint tame
quotient to have order four. One individual image must
therefore have order twelve. If both had that order,
their tame quotient orders and hence their prime-to-$3$
characters would agree. The other image instead has
tame quotient order one or two, and hence total order
three or six. Joint wild order nine is exactly the
asserted independence. Every failure of the unrestricted
$81$ test has now been reduced to (B-RD32.3); the converse
implication about the least index is immediate because
these specified pairs are themselves allowed in B.4.
\end{proof}

\begin{remark}[Exact auxiliary certificate]
The retained program
\path{code/verify_height_two_character_gate.py}
checks the eight cyclotomic reduced norms and traces, all twenty-seven
numerical joint-image possibilities, and the two remaining level-$81$
pairs. Its exact run passed in $0.050076540999725694$ seconds.
The numerical level-$27$ list is an upper list of possibilities; no
existence claim is made for a listed group-theoretic combination.
\end{remark}
